\notechapter{N-20221205}{Counting by Logic, Projection, Inversion, and Generating Series}

\section{Set operations are logical operations written geometrically}

For a finite universe \(U\), the indicator of a set \(A\subseteq U\) is
\[
  \mathbf1_A(x)=
  \begin{cases}
    1,&x\in A,\\
    0,&x\notin A.
  \end{cases}
\]
Then
\[
  |A|=\sum_{x\in U}\mathbf1_A(x).
\]
Union, intersection, and complement become pointwise arithmetic:
\[
  \mathbf1_{A\cap B}=\mathbf1_A\mathbf1_B,
\]
\[
  \mathbf1_{A^c}=1-\mathbf1_A,
\]
\[
  \mathbf1_{A\cup B}
  =\mathbf1_A+\mathbf1_B-\mathbf1_A\mathbf1_B.
\]
De Morgan's law follows immediately:
\[
  1-\mathbf1_{A\cup B}
  =(1-\mathbf1_A)(1-\mathbf1_B).
\]
The set identity
\[
  (A\cup B)^c=A^c\cap B^c
\]
is therefore the logical identity
\[
  \neg(P\lor Q)=(\neg P)\land(\neg Q).
\]

This translation is more than notation.  It allows a set identity to be verified pointwise and then summed to obtain a counting identity.  The same habit will later turn inclusion--exclusion into an algebraic inversion formula.

\section{Addition counts disjoint cases; multiplication counts successive choices}

If a finite collection is partitioned into disjoint cases \(C_i\), then
\[
  \left|\bigsqcup_i C_i\right|
  =\sum_i|C_i|.
\]
If an object is built through successive independent choices, with \(m_i\) options at stage \(i\), then the number of constructions is
\[
  \prod_i m_i.
\]
The first principle counts alternatives; the second counts stages.

Arranging \(r\) distinct objects chosen from \(n\) gives
\[
  n(n-1)\cdots(n-r+1)
  =\frac{n!}{(n-r)!}.
\]
If order is forgotten, every \(r\)-element subset has \(r!\) internal orders, so
\[
  \binom nr
  =\frac{n!}{r!(n-r)!}.
\]
The symmetry
\[
  \binom nr=\binom n{n-r}
\]
comes from the complement bijection: choosing the selected elements is equivalent to choosing the omitted ones.

The pigeonhole principle is an averaging consequence of the same counting.  If \(N\) objects occupy \(r\) boxes, some box contains at least
\[
  \left\lceil\frac Nr\right\rceil
\]
objects; otherwise the total would be smaller than \(N\).

\section{The binomial theorem counts choices inside a product}

Expanding
\[
  (a+b)^n
  =(a+b)\cdots(a+b)
\]
requires choosing either \(a\) or \(b\) from each factor.  A term \(a^{n-r}b^r\) occurs once for every choice of the \(r\) factors that contribute \(b\).  Hence
\[
  \boxed{
  (a+b)^n
  =\sum_{r=0}^n\binom nr a^{n-r}b^r.}
\]
Setting \(a=b=1\) gives
\[
  \sum_{r=0}^n\binom nr=2^n.
\]
This is also the count of all subsets of an \(n\)-element set: each element independently receives the decision “in” or “out.”

Pascal's identity
\[
  \binom nr
  =\binom{n-1}r+\binom{n-1}{r-1}
\]
comes from partitioning \(r\)-subsets according to whether they contain one distinguished element.  The identity is a decomposition of the counted set, not merely a factorial manipulation.

\section{Double counting is counting one incidence set through two projections}

Let \(X\) be a finite set of pairs.  The two coordinate projections partition \(X\) in different ways, and equality of the resulting totals produces an identity.

For example, count pairs
\[
  (S,s),
  \qquad
  S\subseteq[n],
  \quad
  s\in S.
\]
Counting first by \(S\) gives
\[
  |X|=\sum_{k=0}^n k\binom nk.
\]
Counting first by \(s\), choose the distinguished element in \(n\) ways and then choose any subset of the remaining \(n-1\) elements:
\[
  |X|=n2^{n-1}.
\]
Thus
\[
  \boxed{
  \sum_{k=0}^n k\binom nk=n2^{n-1}.}
\]

Vandermonde's identity is another projection count.  Split a population into disjoint groups of sizes \(r\) and \(s\).  Choosing \(n\) people altogether can be done directly in
\[
  \binom{r+s}{n}
\]
ways, or by choosing \(k\) from the first group and \(n-k\) from the second:
\[
  \boxed{
  \sum_k\binom rk\binom s{n-k}
  =\binom{r+s}{n}.}
\]
Naming the incidence object explains the summation index and the range automatically.

\section{Inclusion--exclusion is a pointwise product expansion}

For sets \(A_1,\ldots,A_m\), the indicator of the complement of their union is
\[
  \mathbf1_{(\cup_iA_i)^c}
  =\prod_{i=1}^m(1-\mathbf1_{A_i}).
\]
Expanding and subtracting from one gives
\[
\begin{aligned}
  \mathbf1_{\cup_iA_i}
  ={}&\sum_i\mathbf1_{A_i}
  -\sum_{i<j}\mathbf1_{A_i\cap A_j}\\
  &+\sum_{i<j<k}\mathbf1_{A_i\cap A_j\cap A_k}
  -\cdots.
\end{aligned}
\]
Summing over \(U\) yields the full inclusion--exclusion formula
\[
  \left|\bigcup_{i=1}^mA_i\right|
  =
  \sum_{\varnothing\ne S\subseteq[m]}
  (-1)^{|S|+1}
  \left|\bigcap_{i\in S}A_i\right|.
\]
The alternating signs arise from expanding a product, and every element is checked locally according to how many sets contain it.

As a quick application, the number of surjections from an \(n\)-element set onto an \(m\)-element set is
\[
  \sum_{j=0}^m(-1)^j
  \binom mj(m-j)^n.
\]
Start with all \(m^n\) functions and exclude those missing at least one target value.  Intersections correspond to missing a specified set of \(j\) values, leaving \((m-j)^n\) functions.

\section{Boolean Möbius inversion explains the alternating signs structurally}

Let functions \(f,g\) be defined on subsets of \([m]\), with
\[
  g(S)=\sum_{T\subseteq S}f(T).
\]
This cumulative transform is the zeta transform of the Boolean lattice.  Möbius inversion recovers the exact data:
\[
  \boxed{
  f(S)=
  \sum_{T\subseteq S}
  (-1)^{|S|-|T|}g(T).}
\]
The Boolean Möbius function is therefore
\[
  \mu(T,S)=(-1)^{|S|-|T|}
  \qquad(T\subseteq S).
\]

For \(m=3\), the top value is
\[
\begin{aligned}
  f(123)={}&g(123)
  -g(12)-g(13)-g(23)\\
  &+g(1)+g(2)+g(3)-g(\varnothing).
\end{aligned}
\]
The familiar inclusion--exclusion signs are exactly the coefficients of the inverse transform.

To verify the formula, substitute the definition of \(g\).  The coefficient of a fixed \(f(R)\) is
\[
  \sum_{R\subseteq T\subseteq S}
  (-1)^{|S|-|T|}
  =(1-1)^{|S|-|R|},
\]
which is \(1\) when \(R=S\) and \(0\) otherwise.  Inversion works because all unwanted lower-level contributions cancel.

\section{Ordinary generating functions turn decomposition into multiplication}

For a sequence \((a_n)_{n\ge0}\), its ordinary generating function is
\[
  A(x)=\sum_{n\ge0}a_nx^n.
\]
The notation
\[
  [x^n]A(x)=a_n
\]
extracts a coefficient.  If an object of size \(n\) splits into an object of size \(k\) and one of size \(n-k\), the number of pairs is the convolution
\[
  c_n=\sum_{k=0}^n a_kb_{n-k},
\]
and
\[
  C(x)=A(x)B(x).
\]
The multiplication principle has become multiplication of series.

Vandermonde's identity follows algebraically from
\[
  (1+x)^r(1+x)^s=(1+x)^{r+s}.
\]
Extracting \([x^n]\) from the left gives
\[
  \sum_k\binom rk\binom s{n-k},
\]
while the right gives \(\binom{r+s}n\).  The same identity now has both an incidence-set proof and a coefficient proof.

\section{Generating functions solve recurrences and coin-change problems}

Let \(F_0=0\), \(F_1=1\), and
\[
  F_n=F_{n-1}+F_{n-2}.
\]
Set
\[
  F(x)=\sum_{n\ge0}F_nx^n.
\]
Multiplying the recurrence by \(x^n\) and summing for \(n\ge2\) gives
\[
  F(x)-x=xF(x)+x^2F(x).
\]
Therefore
\[
  \boxed{
  F(x)=\frac{x}{1-x-x^2}.}
\]
Factoring the denominator recovers Binet's formula; alternatively, the rational function generates the sequence without solving each recurrence anew.

For unlimited coins of values \(1,2,5\), choosing \(a,b,c\ge0\) coins contributes value
\[
  a+2b+5c.
\]
The generating function is
\[
  \frac1{(1-x)(1-x^2)(1-x^5)}.
\]
The coefficient of \(x^n\) counts solutions of
\[
  a+2b+5c=n.
\]
For \(n=5\), the four representations are
\[
  1+1+1+1+1,
  \quad
  2+1+1+1,
  \quad
  2+2+1,
  \quad
  5,
\]
so
\[
  [x^5]\frac1{(1-x)(1-x^2)(1-x^5)}=4.
\]
A product of geometric series encodes independent multiplicities of allowed parts.

\section{Averaging and random choice can prove existence}

The pigeonhole principle says that some box is at least as full as the average.  The probabilistic method applies the same idea to a random object: if the expected value of a statistic is \(M\), some object has value at least \(M\).

Let a graph have \(m\) edges.  Place each vertex independently into one of two parts with probability \(1/2\).  An edge crosses the cut exactly when its endpoints land in different parts, which occurs with probability \(1/2\).  If \(X\) is the number of crossing edges, linearity of expectation gives
\[
  \mathbb E[X]
  =\sum_{e}\mathbb P(e\text{ crosses})
  =\frac m2.
\]
Therefore at least one cut satisfies
\[
  X\ge\frac m2.
\]
No independence between edge indicators is required for linearity of expectation.

The argument is constructive in logic even when it does not exhibit the cut: if every cut had fewer than \(m/2\) crossing edges, the average could not equal \(m/2\).

\section{Stirling's formula reveals the scale of the central binomial coefficient}

The largest coefficient in the \(n\)-th row of Pascal's triangle lies near \(n/2\).  For even \(n=2m\), Stirling's approximation
\[
  m!\sim\sqrt{2\pi m}\left(\frac me\right)^m
\]
gives
\[
\begin{aligned}
  \binom{2m}{m}
  &=\frac{(2m)!}{(m!)^2}\\
  &\sim
  \frac{\sqrt{4\pi m}(2m/e)^{2m}}
       {2\pi m(m/e)^{2m}}\\
  &=\frac{4^m}{\sqrt{\pi m}}.
\end{aligned}
\]
Equivalently,
\[
  \boxed{
  \binom n{\lfloor n/2\rfloor}
  \sim
  2^n\sqrt{\frac{2}{\pi n}}.}
\]
Thus the largest layer of the Boolean lattice contains about \(2^n/\sqrt n\) subsets.  The exact identity
\[
  \sum_k\binom nk=2^n
\]
counts the whole lattice; the asymptotic formula describes how that mass concentrates around the middle ranks.

\section{Species record how labelled structures move under relabelling}

A combinatorial species \(F\) assigns to each finite label set \(U\) a set \(F[U]\) of structures, and every bijection
\[
  \sigma:U\to V
\]
transports structures through a bijection
\[
  F[\sigma]:F[U]\to F[V]
\]
compatible with composition.  The definition remembers not only how many structures exist, but how they behave under relabelling.

The exponential generating function is
\[
  F(x)=
  \sum_{n\ge0}
  |F[n]|\frac{x^n}{n!}.
\]
For the species \(\mathrm{SET}\), there is one set structure on every label set, so
\[
  \mathrm{SET}(x)
  =\sum_{n\ge0}\frac{x^n}{n!}
  =e^x.
\]
For the species \(\mathrm{SEQ}\), a label set of size \(n\) has \(n!\) linear orders, so
\[
  \mathrm{SEQ}(x)
  =\sum_{n\ge0}n!\frac{x^n}{n!}
  =\frac1{1-x}.
\]
For \(n\ge1\), a labelled cycle has \((n-1)!\) cyclic orders, hence
\[
  \mathrm{CYC}(x)
  =\sum_{n\ge1}(n-1)!\frac{x^n}{n!}
  =\sum_{n\ge1}\frac{x^n}{n}
  =-\log(1-x).
\]

Every permutation decomposes uniquely into an unordered set of disjoint cycles.  Species composition therefore gives
\[
  \mathrm{PERM}=\mathrm{SET}(\mathrm{CYC}).
\]
At the level of exponential generating functions,
\[
\begin{aligned}
  \mathrm{PERM}(x)
  &=\exp(\mathrm{CYC}(x))\\
  &=\exp\bigl(-\log(1-x)\bigr)\\
  &=\frac1{1-x}.
\end{aligned}
\]
The coefficient of \(x^n/n!\) is \(n!\), recovering the number of permutations.  Addition, product, and composition of labelled structures have become algebraic operations on generating series because relabelling symmetry is built into the objects from the beginning.